\section{Modelo de dominio}

\begin{figure}[htbp]
\centering
\includegraphics[width=\textwidth,height=0.93\textheight,keepaspectratio]{"../../architecture/SDK/Domain Class Diagram.pdf"}
\caption{Diagrama de clases del dominio del SDK}
\label{fig:domain_classes_SDK}
\end{figure}

\FloatBarrier
\renewcommand{\arraystretch}{1.3}
\begin{longtable}{|>{\raggedright\arraybackslash}p{0.25\textwidth}|p{\dimexpr\textwidth-0.25\textwidth-4\tabcolsep-3\arrayrulewidth\relax}|}
\caption{Conceptos del SDK y su explicación}
\label{tab:domain_concepts_SDK}\\
\hline
\textbf{Concepto} & \textbf{Explicación} \\ \hline
\endfirsthead
\hline
\textbf{Concepto} & \textbf{Explicación} \\ \hline
\endhead
\hline
\endfoot
\endlastfoot
Transaction (Entidad) & Representa una transacción de pago tal como la reporta la pasarela: agrupa el identificador nativo, la referencia del comercio, el monto, la moneda, los datos del pagador, el estado normalizado y el nativo y, cuando lo hay, el motivo de rechazo. Es inmutable (punto 3 del registro de arquitectura): un estado nuevo produce otra instancia y la anterior queda igual. Es la única entidad del modelo, es decir, la única clase que se distingue por su identidad y no por sus valores. \\ \hline
Amount (Objeto de Valor) & Encapsula el monto en unidades monetarias con hasta dos decimales — la representación canónica del SDK. Cada adaptador es responsable de convertirla al formato nativo de su pasarela (centavos, objeto descompuesto, etc.). \\ \hline
Currency (Objeto de Valor) & Encapsula el código ISO 4217 de la moneda (por defecto "COP") y valida que tenga exactamente tres letras mayúsculas. \\ \hline
OrderReference (Objeto de Valor) & Es la referencia que el comercio asigna a la transacción para correlacionarla con sus propios sistemas. El SDK valida que no esté vacía, pero nunca la genera ni garantiza su unicidad — eso es responsabilidad del comercio. \\ \hline
Payer (Objeto de Valor) & Agrupa los datos de quien paga (email, nombre, tipo y número de documento, teléfono). El único campo obligatorio es el email, por ser el único dato que las cuatro pasarelas exigen en común. \\ \hline
ReturnUrlConfig (Objeto de Valor) & Configura a dónde se redirige al usuario después de pagar. Solo aplica en integraciones hosted/redirect. Contempla una URL única o URLs diferenciadas por resultado (success/failure/pending), porque Mercado Pago es la única pasarela con esa granularidad. \\ \hline
GatewayTransactionId (Objeto de Valor) & El identificador que la pasarela genera internamente para la transacción — distinto de OrderReference. Se usa para consultas de estado, reembolsos y conciliación contra la pasarela. \\ \hline
RejectionReason (Objeto de Valor) & Explica por qué se rechazó una transacción, combinando el código nativo de la pasarela (rejectionCode) con una categoría semántica normalizada (rejectionCategory) que el comercio puede usar para decidir su lógica de negocio. \\ \hline
WebhookVerifier (Servicio de Dominio) & Verifica que un evento de webhook realmente venga de la pasarela y no haya sido alterado ni falsificado. Se modela como servicio, no como entidad ni objeto de valor, porque no tiene datos propios que persistir — es comportamiento puro. \\ \hline
KitPagosError (Error de Dominio) & El único tipo de error que el SDK expone hacia afuera, sin importar qué pasarela lo originó. Extiende el Error nativo y añade \texttt{code: KitPagosErrorCode}, \texttt{gateway: Gateway}, \texttt{originalPayload: unknown} y una causa opcional, de modo que el comercio pueda reaccionar al fallo sin interpretar excepciones nativas ni payloads crudos. \\ \hline
TransactionStatus (Enumeración) & Los seis estados posibles de una transacción, ya normalizados entre las cuatro pasarelas: APPROVED, DECLINED, PENDING, EXPIRED, VOIDED, ERROR. Es el único vocabulario de estado que el código del comercio debe conocer. \\ \hline
RejectionCategory (Enumeración) & Clasifica el motivo de un rechazo en ocho categorías de negocio (fondos insuficientes, datos inválidos, tarjeta restringida, fraude, error de red, autenticación fallida, pago duplicado, desconocido), para que la lógica de reintento no dependa de interpretar strings distintos por pasarela. \\ \hline
KitPagosErrorCode (Enumeración) & Códigos de error técnico normalizados (timeout, credenciales inválidas, límite de solicitudes, operación no soportada, etc.), independientes de qué pasarela los causó. Son trece, y \texttt{KitPagosError} lleva uno de ellos en \texttt{code}. \\ \hline
Gateway (Enumeración) & Identifica cuál de las cuatro pasarelas está involucrada (WOMPI, RAPYD, MERCADOPAGO, KUSHKI). Se usa tanto en GatewayTransactionId como en KitPagosError para saber de dónde vino cada dato. \\ \hline
PaymentMethod (Objeto de Valor) & Describe el medio de pago con el que se cobra: \texttt{CARD} o \texttt{PSE}. Se construye con las fábricas \texttt{PaymentMethod.card(cardToken?, \{ installments? \})} y \texttt{PaymentMethod.pse(\{ bankCode, payerKind? \})}; \texttt{requiresPayerDocument()} indica si el flujo exige documento del pagador, como ocurre en PSE. \\ \hline
PaymentResult y PendingRedirect (Objeto de Valor) & Unión discriminada por \texttt{outcome} que retorna \texttt{createPayment()}: \texttt{TRANSACTION} cuando el cobro queda resuelto en una \texttt{Transaction}, o \texttt{REDIRECT\_REQUIRED} cuando el comprador debe completar el pago en una URL externa. \texttt{PendingRedirect} agrupa \texttt{redirectUrl}, \texttt{gatewayTransactionId} y \texttt{rawStatus}. \\ \hline
TaxBreakdown (Objeto de Valor) & Descompone un monto en \texttt{subtotalIva}, \texttt{subtotalIva0}, \texttt{iva} e \texttt{ice}, y expone su total con \texttt{getTotal()}. Lo exige Kushki, que recibe el monto desglosado en lugar de un valor único; las demás pasarelas usan \texttt{TaxBreakdown.exempt()}. \\ \hline
PseBank y PseBankCode (Objeto de Valor) & \texttt{PseBank} representa una entidad financiera habilitada para PSE (\texttt{code}, \texttt{name} y \texttt{achCode} opcional) y es lo que devuelve \texttt{getPseBanks()}. \texttt{PseBankCode} enumera los códigos de compensación de ACH Colombia y permite traducirlos al código que exige cada proveedor. \\ \hline
Environment (Objeto de Valor) & Define el ambiente de ejecución del SDK: \texttt{simulator}, \texttt{sandbox} o \texttt{production}. A partir de él, \texttt{SdkConfigurator} resuelve el destino desde el catálogo cerrado de URLs, sin que el comercio escriba dominios de proveedores ni pueda inyectar una URL arbitraria. \\ \hline
Credentials (Objeto de Valor) & Encapsula las llaves de una pasarela: \texttt{publicKey} y \texttt{privateKey}, más los opcionales \texttt{integritySecret} (Wompi), \texttt{webhookSecret} (Wompi, Mercado Pago y Kushki) y \texttt{webhookUrl} (Rapyd, que firma la URL del receptor). \texttt{SdkConfigurator} mantiene un mapa \texttt{Gateway} $\to$ \texttt{Credentials} y cada adaptador las usa sin propagarlas a logs, errores ni respuestas (RF-08). \\ \hline
WebhookEvent (Objeto de Valor) & Representa el evento normalizado que produce un webhook válido: \texttt{eventType}, \texttt{gatewayTransactionId}, \texttt{newStatus}, \texttt{gateway} y, cuando la firma del proveedor cubre el cuerpo, \texttt{rawStatus} y \texttt{orderReference}. Es lo que retorna \texttt{validateWebhook()}, en lugar de un booleano. \\ \hline
\end{longtable}

\begin{figure}[htbp]
\centering
\includegraphics[width=\textwidth,height=0.93\textheight,keepaspectratio]{"../../architecture/Simulador API/diagrama-clases-dominio-simulacion.pdf"}
\caption{Diagrama de clases del dominio de la API de Simulación}
\label{fig:domain_classes_API}
\end{figure}

\begin{table}[htbp]
\centering
\caption{Conceptos Adicionales del Dominio (Simulación)}
\label{tab:domain_concepts_API}
\renewcommand{\arraystretch}{1.3}
\begin{tabularx}{\textwidth}{|l|X|}
\hline
\textbf{Concepto} & \textbf{Explicación} \\ \hline
SimulationScenario (Enumeración) & Los escenarios controlables que un desarrollador o tester puede solicitar al simulador, normalizados en inglés para ser consistentes con el resto del vocabulario del dominio: \texttt{APPROVED}, \texttt{DECLINED}, \texttt{PENDING}, \texttt{EXPIRED}, \texttt{DUPLICATE\_PAYMENT}, \texttt{RATE\_LIMIT}, \texttt{SERVER\_ERROR}, \texttt{BAD\_GATEWAY}, \texttt{SERVICE\_UNAVAILABLE}, \texttt{TIMEOUT}, \texttt{NETWORK\_ERROR}, \texttt{FLAPPING}, \texttt{SLOW}, \texttt{MALFORMED\_JSON}, \texttt{MALFORMED\_BODY}, \texttt{HTML\_ERROR} e \texttt{INVALID\_CREDENTIALS}, además de las fallas de consulta \texttt{QUERY\_SERVER\_ERROR}, \texttt{QUERY\_FLAPPING} y \texttt{QUERY\_SLOW}. El escenario se resuelve en la creación con la precedencia cabecera, dato de prueba, monto reservado y omisión; las consultas no lo alteran. \\ \hline
WebhookTriggerRequest (Objeto de Valor) & Forma de la solicitud que recibe \texttt{POST /v1/sim/webhooks/trigger}: solo \texttt{gateway} y \texttt{transactionId}. La URL destino no viaja en el cuerpo, sino que sale de \texttt{SIMULATOR\_WEBHOOK\_TARGET\_URL}; cualquier otra clave se rechaza con \texttt{INVALID\_REQUEST}. El simulador también emite webhooks de forma automática cuando cambia el estado de un pago, si \texttt{SIMULATOR\_WEBHOOK\_AUTO} está habilitado (RF-12). \\ \hline
\end{tabularx}
\end{table}


